% Keep the temporary build directory out of the PDF trailer ID.
\ifdefined\pdfvariable\pdfvariable suppressoptionalinfo 512\relax\fi
\documentclass{ltjarticle}
\title{格子窓の試し刷り}
\author{紙窓模型工房}
\date{}
\begin{document}
\maketitle
\begin{abstract}
窓枠を紙に写し、折り返した帯が開口を隠さないか調べた。
\end{abstract}
\tableofcontents
\section{採寸}
\subsection{帯の向き}
白い面を上にして窓枠を置く。\footnote{台紙の折り目を基準にする。}
\paragraph{目印}
長辺に一本の筋を引く。
\begin{itemize}\item 桟の幅を写す。\item 短辺の切り口をそろえる。\end{itemize}
\newpage
\section{試験結果}
図\ref{fig:bar}の桟を表\ref{tab:width}の順に重ねる。
\begin{figure}[htbp]
\centering\rule{30mm}{3mm}
\caption{窓枠の桟}\label{fig:bar}
\end{figure}
\begin{table}[htbp]
\centering\caption{帯の幅}\label{tab:width}
\begin{tabular}{lr}位置 & 幅（ミリ）\\内側 & 3\\外側 & 5\end{tabular}
\end{table}
\begin{enumerate}\item 白紙を窓に当てる。\item 影の幅を台帳に写す\cite{ledger}。\end{enumerate}
\begin{thebibliography}{9}
\bibitem{ledger} 紙窓模型工房『採光台帳』。
\end{thebibliography}
\end{document}
